\section{Method}

Our framework builds upon VGGT and extends it from static 3D perception to dynamic 4D reconstruction. 
The key idea is to establish a unified geometric representation, namely DPMs, as the core of temporal modeling. 
Based on this formulation, we introduce temporal reasoning via a MTA module, predict future geometry via a Future Point Head (FPH), and further refine dynamic geometry through a DGSHead. 
An overview of the proposed architecture is shown in Fig.~\ref{fig:pipeline}.

\subsection{Dynamic Point Map and Task Formulation}

% We denote the camera/view index by $v \in \{1,\dots,N_v\}$ and the discrete frame index by $t \in \{1,\dots,\tau\}$, where $N_v$ is the number of camera views and $\tau$ is the clip length. 
% Unless otherwise specified, temporal prediction is defined on frame pairs $(v,t)$ and $(v,t+\delta)$ from the same camera stream, where $\delta$ denotes a discrete temporal offset.
% Modeling point-wise motion is expected to facilitate the learning of dynamic information within the network. 
% Previous works~\cite{sucar2025dynamic} define dynamic point maps on top of DUSt3R~\cite{wang2024dust3r}. 
% In the static case, the point map at time $t$ for view $v$ is formulated as
% \begin{equation}
% P_{v,t} = \pi^{-1}(I_{v,t}; K_{v,t}, E_{v,t}) \in \mathbb{R}^{3 \times H \times W},
% \end{equation}
% where $\pi^{-1}(\cdot)$ denotes the back-projection from image coordinates to 3D space using the camera intrinsics $K_{v,t}$ and extrinsics $E_{v,t}$. 
% Each pixel $(x,y)$ in $I_{v,t}$ corresponds to a 3D point $\mathbf{p}_{v,t}(x,y)$ defined in the camera coordinate system at time $t$.
% To model dynamics, prior DPM formulations~\cite{sucar2025dynamic} introduce a temporal extension under a fixed reference viewpoint and time condition. 
% Instead of defining each frame in its own coordinate system, all frames are expressed in a shared reference frame:
% \begin{equation}
% P_{v,t}^{(\mathrm{ref})} = \mathcal{T}_{(v,t)\rightarrow \mathrm{ref}}
% \big( \pi^{-1}(I_{v,t}; K_{v,t}, E_{v,t}) \big),
% \end{equation}
% where $\mathcal{T}_{(v,t)\rightarrow \mathrm{ref}}$ denotes the transformation that aligns the 3D points from frame $(v,t)$ to the reference frame. 
% Temporal correspondence can then be modeled as
% \begin{equation}
% \Delta P_{v,t}^{(\mathrm{ref})} = P_{v,t+\delta}^{(\mathrm{ref})} - P_{v,t}^{(\mathrm{ref})},
% \end{equation}
% which represents point-wise 3D motion in a temporally consistent coordinate space.


% In our formulation, the VGGT backbone inherently predicts point maps in a learned canonical coordinate system. 
% Therefore, instead of explicitly constructing dynamic point maps from externally provided camera transformations, we directly regress point maps through a feed-forward model. 
% Given a multi-view input clip $\{I_{v,t}\}$, the network predicts
% \begin{equation}
% \hat{P}_{v,t}, \hat{P}_{v,t+\delta} = f_{\theta}\!\left(\{I_{v,t}\}\right)\Big|_{(v,t),(v,t+\delta)},
% \end{equation}
% where $f_{\theta}$ denotes the proposed DynamicVGGT. 
% Both predicted point maps are represented in a unified canonical frame, enabling the model to learn point-wise motion implicitly through
% $
% \Delta \hat{P}_{v,t} = \hat{P}_{v,t+\delta} - \hat{P}_{v,t}.
% $
% This formulation avoids explicitly relying on externally specified frame-to-reference transformations in the dynamic formulation, while still preserving the geometric prior of the original VGGT backbone. 

% It therefore provides a unified basis for dynamic task design, as illustrated in Fig.~\ref{fig:dynamic_task}.
% Specifically, we introduce two complementary tasks built upon this formulation: 
% (i) the FPH, which learns implicit motion by enforcing inter-frame point consistency, and 
% (ii) the DGSHead, which explicitly refines dynamic geometry through scene-flow-supervised Gaussian optimization.

We denote the camera index by $v \in \{1,\dots,N_v\}$ and the frame index by $t \in \{1,\dots,\tau\}$. 
Temporal modeling is defined on frame pairs $(v,t)$ and $(v,t+\delta)$ from the same camera stream.
Following prior work~\cite{sucar2025dynamic}, a static point map is defined as
\begin{equation}
P_{v,t} = \pi^{-1}(I_{v,t}; K_{v,t}, E_{v,t}) \in \mathbb{R}^{3 \times H \times W}.
\end{equation}

To model dynamics, previous DPM formulations align all frames into a shared reference frame:
\begin{equation}
P_{v,t}^{(\mathrm{ref})} =
\mathcal{T}_{(v,t)\rightarrow \mathrm{ref}}
\big(\pi^{-1}(I_{v,t};K_{v,t},E_{v,t})\big),
\end{equation}

so that temporal motion is expressed as
\begin{equation}
\Delta P_{v,t}^{(\mathrm{ref})} =
P_{v,t+\delta}^{(\mathrm{ref})}-P_{v,t}^{(\mathrm{ref})}.
\end{equation}

In contrast, VGGT directly predicts point maps in a learned canonical frame. 
Given a multi-view clip $\{I_{v,t}\}$, our model predicts
\begin{equation}
\hat{P}_{v,t},\hat{P}_{v,t+\delta}=
f_\theta(\{I_{v,t}\})|_{(v,t),(v,t+\delta)},
\end{equation}

which enables implicit motion learning through
$
\Delta\hat{P}_{v,t}=\hat{P}_{v,t+\delta}-\hat{P}_{v,t}.
$
This formulation avoids explicitly relying on externally specified frame-to-reference transformations in the dynamic formulation, while still preserving the geometric prior of the original VGGT backbone. 

It therefore provides a unified basis for dynamic task design, as illustrated in Fig.~\ref{fig:dynamic_task}.
Specifically, we introduce two complementary tasks built upon this formulation: 
(i) the FPH, which learns implicit motion by enforcing inter-frame point consistency, and 
(ii) the DGSHead, which explicitly refines dynamic geometry through scene-flow-supervised Gaussian optimization.

\begin{figure}[t]
\centering
  \includegraphics[width=0.9\columnwidth]{figs/dynamic_task_design.pdf}
  \caption{\textbf{Dynamic task formulation.} 
  We formulate dynamic point maps by designing two complementary tasks that model point-wise motion over time. 
  The Future Point Head learns implicit motion through inter-frame point consistency, while the Dynamic 3D Gaussian Splatting Head provides explicit motion supervision via scene flow to refine dynamic geometry.}
  \label{fig:dynamic_task}
\end{figure}

\subsection{Motion-aware Temporal Attention}

Although the DPM formulation captures temporal variations at the geometric level, relying solely on point-wise displacement is still insufficient for accurate dynamic reconstruction. 
Previous work, such as StreamVGGT~\cite{zhuo2025streaming}, also introduces temporal attention to enhance temporal modeling. 
However, its sequential stacking of AA blocks tends to cause unstable training and degraded performance in the early stages.
To address this limitation, we propose a Motion-aware Temporal Attention (MTA) module that explicitly models motion cues at the feature level and integrates temporal reasoning into the VGGT backbone. 
The key idea is to introduce learnable motion tokens that dynamically encode inter-frame motion information, guiding temporal attention to focus on motion-consistent regions.

Given the aggregated token features $\tilde{F}_{v,t} = [F_{v,t}^{c}; F_{v,t}^{p}]$ produced by the AA blocks, we remove the camera token $F_{v,t}^{c}$ and concatenate the patch tokens $F_{v,t}^{p}$ with learnable motion tokens to form the MTA input. 
For each MTA layer $l$, temporal correlations are computed in parallel across all frames along the temporal dimension $\tau$. 
The input to the $l$-th MTA layer is defined as
\begin{equation}
F_{m,v,t}^{(l)} = 
\begin{cases}
\mathrm{Concat}\!\left(M_{v,t}^{(l)},\, F_{v,t}^{p(l)}\right), & l = 1,\\
\mathrm{Concat}\!\left(M_{v,t}^{(l)},\, F_{v,t}^{p(l)} + F_{v,t}^{p(l-1)}\right), & l > 1,
\end{cases}
\end{equation}
where $M_{v,t}^{(l)}$ denotes the motion tokens and $F_{v,t}^{p(l)}$ denotes the spatial patch tokens from the AA branch at layer $l$.
% 中文：这里把所有 token 都补上了 view index。
% 中文：同时把时间长度 T 改成 \tau，避免和 transformation T 混淆。
Temporal attention is computed independently for each patch position and each view:
\begin{equation}
A_{t,t'}^{(l)} = 
\mathrm{Softmax}\!\left(
\frac{Q_{t}^{\mathrm{attn},(l)} \left(K_{t'}^{\mathrm{attn},(l)}\right)^\top}{\sqrt{d}}
+ B_{t,t'}^{\mathrm{time}}
\right),
\end{equation}
\begin{equation}
\tilde{F}_{m,v,t}^{(l)} = \sum_{t'=1}^{\tau} A_{t,t'}^{(l)} V_{t'}^{\mathrm{attn},(l)},
\end{equation}
where $t,t' \in \{1,\dots,\tau\}$ denote discrete frame indices within the sampled clip, and $B_{t,t'}^{\mathrm{time}}$ denotes the temporal positional bias implemented using rotary position embeddings.
The updated temporal features are then processed by layer normalization and an MLP with residual connections:
\begin{equation}
F_{m,v,t}^{(l+1)} = \mathrm{MLP}^{(l)}\!\left(\mathrm{LayerNorm}\!\left(\tilde{F}_{m,v,t}^{(l)}\right)\right) + F_{m,v,t}^{(l)}.
\end{equation}

After the final MTA layer, we denote the temporally enhanced feature for view $v$ at time $t$ by
\begin{equation}
TA_{v,t} = F_{m,v,t}^{(L)},
\end{equation}
where $L$ is the number of MTA layers. 
The resulting feature $TA_{v,t}$ is used by both the Future Point Head and the Dynamic 3DGS Head.
This formulation enables simultaneous message passing across temporal spans and enhances the model’s ability to capture motion continuity and temporally coherent geometry.

\subsection{Future Point Prediction}

Building upon the unified DPM representation and the feature-level temporal modeling provided by MTA, 
we further introduce a Future Point Head (FPH) to explicitly learn point-wise motion dynamics. 
Specifically, the FPH predicts the 3D point map of a future frame from the temporally enhanced feature at the current timestep, enabling the network to learn short-term motion continuity in a self-supervised manner.
Given the temporally enhanced feature $TA_{v,t}$ produced by the MTA module, the FPH predicts the point map of the same camera stream at a future timestep:
\begin{equation}
\hat{P}_{v,t+\delta}^{\,\mathrm{fut}} = \mathrm{DPT}_{p}(TA_{v,t}),
\end{equation}
where $\mathrm{DPT}_{p}(\cdot)$ denotes a DPT head~\cite{ranftl2021vision} for future point regression.
To further encourage physically plausible point motion trajectories, we introduce a temporal consistency regularization:
\begin{equation}
\mathcal{L}_{\mathrm{temp}}=
\frac{1}{|\mathcal{N}|}\sum_{i\in \mathcal{N}}
\left\|
\left(\mathbf{p}_{v,t+\delta}^{(i)}-\mathbf{p}_{v,t}^{(i)}\right)
-
\left(\hat{\mathbf{p}}_{v,t+\delta}^{(i)}-\hat{\mathbf{p}}_{v,t}^{(i)}\right)
\right\|_{1},
\label{eq:temporal_consistency}
\end{equation}
where $\mathbf{p}_{v,t}^{(i)}$ and $\hat{\mathbf{p}}_{v,t}^{(i)}$ denote the ground-truth and predicted 3D coordinates of the $i$-th valid point at time $t$, respectively, and $\mathcal{N}$ denotes the set of valid points.
The displacement field
$
\Delta \mathbf{p}_{v,t}^{(i)} = \mathbf{p}_{v,t+\delta}^{(i)}-\mathbf{p}_{v,t}^{(i)}
$
acts as a coarse motion representation that encourages the network to learn inter-frame point displacement within the shared DPM coordinate space. 
We emphasize that $\mathcal{L}_{\mathrm{temp}}$ supervises motion implicitly at the point-map level, which complements the explicit motion supervision introduced later in the Dynamic 3DGS Head.

\subsection{Dynamic 3D Gaussian Splatting Head}

To further model dynamic scenes at the primitive level, we introduce a Dynamic 3D Gaussian Splatting (3DGS) Head as a downstream dynamic reconstruction task.

This module takes both the temporally enhanced features $TA_{v,t}$ from MTA and the RGB appearance cues from the input images as inputs, and converts them into time-varying 3D Gaussian primitives that jointly model geometry, appearance, and motion. 
We observe that freezing the AA blocks causes the pretrained VGGT backbone to overemphasize geometric reasoning while weakening appearance cues, which degrades Gaussian rendering quality. 
To compensate for this issue, we fuse image features extracted by $\mathrm{Conv}(\cdot)$ with geometric features obtained from $TA_{v,t}$:
\begin{equation}
F_{v,t}^{\mathrm{app}} = \mathrm{Conv}(I_{v,t}),
\end{equation}
\begin{equation}
F_{g,v,t},\, D_{g,v,t} = \mathrm{DPT}_{g}(TA_{v,t}),
\end{equation}
\begin{equation}
G_{v,t} = F_{v,t}^{\mathrm{app}} + F_{g,v,t},
\end{equation}
where $F_{g,v,t}$ denotes the Gaussian features used to initialize Gaussian primitives, and $D_{g,v,t}$ denotes the predicted Gaussian depth.
At each timestep, the predicted Gaussian depth $D_{g,v,t}$ together with the camera predictions from the retained VGGT camera branch are used to reconstruct a point map $P_{v,t}^{g}$, which initializes the Gaussian centers $\mu_i$. 
Each Gaussian primitive is parameterized as
$
\{\mu_i, \sigma_i, r_i, c_i, \nu_i\},
$
where $\mu_i$ denotes the center, $\sigma_i$ the scale, $r_i$ the rotation, $c_i$ the color, and $\nu_i$ the velocity vector.

We further leverage the $M$ learnable motion tokens introduced in MTA to decode a set of velocity bases $\nu_b \in \mathbb{R}^{3}$, forming a shared dynamic representation for Gaussian motion. 
To describe temporal evolution, we assume constant velocity within a short clip:
\begin{equation}
\mu_{i,t+\delta} = \mu_{i,t} + \delta \cdot \nu_{i,t}.
\end{equation}

During training, scene-flow supervision is applied to encourage each Gaussian primitive to carry a physically meaningful velocity attribute. 
Unlike $\mathcal{L}_{\mathrm{temp}}$, which constrains coarse inter-frame displacement at the point-map level, the Gaussian motion supervision explicitly regularizes motion in the dynamic Gaussian space and therefore provides complementary supervision for dynamic reconstruction.

% 中文：这句再次明确 L_temp 和 L_flow 分工不同，不是重复监督。

\subsection{Training Objective}

To enable the model to learn dynamic geometry progressively, we adopt a two-stage training strategy that follows a curriculum-style paradigm from synthetic to real-world data. 
In the first stage, the model is trained on high-quality synthetic autonomous driving datasets, where dense geometry and reliable motion cues are available. 
The stage-1 training objective combines static and temporal supervision:
\begin{equation}
\mathcal{L}_{\mathrm{stage1}} =
\mathcal{L}_{\mathrm{cam}}
+ \mathcal{L}_{\mathrm{depth}}
+ \mathcal{L}_{\mathrm{point}}^{(t)}
+ \mathcal{L}_{\mathrm{point}}^{(t+\delta)}
+ \lambda_{\mathrm{temp}} \mathcal{L}_{\mathrm{temp}}.
\end{equation}

The camera loss supervises the predicted camera parameters $\hat{g}$:
$
\mathcal{L}_{\mathrm{cam}} = \sum_{i=1}^{N} \left\| \hat{g}_i - g_i \right\|_{\epsilon},
$
where $\hat{g}_i$ and $g_i$ denote the predicted and ground-truth camera parameters, respectively, and $\|\cdot\|_{\epsilon}$ denotes the Huber loss. 
The depth loss $\mathcal{L}_{\mathrm{depth}}$ and point-map losses $\mathcal{L}_{\mathrm{point}}^{(t)}$ and $\mathcal{L}_{\mathrm{point}}^{(t+\delta)}$ follow VGGT~\cite{wang2025vggt}. 
This stage focuses on learning temporal consistency by predicting future point maps, allowing the network to capture short-term motion while preserving the geometric priors of the pretrained backbone.

In the second stage, we fine-tune the model on real driving datasets using the Dynamic 3DGS objective. 
The overall stage-2 objective is defined as
\begin{equation}
\mathcal{L}_{\mathrm{stage2}} = \mathcal{L}_{\mathrm{stage1}} + \mathcal{L}_{\mathrm{3DGS}},
\end{equation}
where
\begin{equation}
\mathcal{L}_{\mathrm{3DGS}} =
\mathcal{L}_{\mathrm{rgb}}
+ \lambda_{\mathrm{gs}} \mathcal{L}_{\mathrm{gsdepth}}
+ \lambda_{\mathrm{dist}} \mathcal{L}_{\mathrm{distill}}
+ \lambda_{\mathrm{flow}} \mathcal{L}_{\mathrm{flow}}.
\end{equation}

\begin{figure}[t]
\centering
  \includegraphics[width=0.9\columnwidth]{figs/point_comparison_v2.pdf}
  \caption{\textbf{Depth and Point Maps Comparison}. The sparsity of LiDAR point clouds degrades the results, leading to less smooth depth maps and rougher point clouds.}
  \label{fig:depth_lidar}
\end{figure}

For image reconstruction, we use
$
\mathcal{L}_{\mathrm{rgb}} = \mathrm{MSE}(I_{v,t}, \hat{I}_{v,t}),
$
where $I_{v,t}$ is the ground-truth RGB image and $\hat{I}_{v,t}$ is the image rendered from the predicted dynamic Gaussian primitives.

We observe that directly using sparse LiDAR point clouds as supervision on real autonomous driving datasets leads to severe performance degradation due to their limited density and uneven spatial distribution. 
To alleviate this issue, we introduce a depth distillation strategy. 
Specifically, the depth predicted by the stage-1 point-map branch serves as the teacher signal, while the Gaussian depth branch predicts the student depth:
$
\mathcal{L}_{\mathrm{distill}} = \left\| D_{g,v,t} - \mathrm{sg}\!\left(D^{\mathrm{pm}}_{v,t}\right) \right\|_1,
$
where $D^{\mathrm{pm}}_{v,t}$ denotes the depth predicted by the stage-1 geometry branch and $\mathrm{sg}(\cdot)$ denotes the stop-gradient operator. 
In addition, $\mathcal{L}_{\mathrm{gsdepth}}$ is a standard $L_1$ loss, whose supervision is provided by the pretrained model~\cite{wang2025moge}. 
This strategy mitigates the noise caused by point-cloud sparsity and stabilizes Gaussian optimization, as shown in Fig.~\ref{fig:depth_lidar}.

\begin{table*}[t]
\centering
\caption[Valori medi]{\textbf{Point Map Reconstruction on KITTI and Waymo(val).} KITTI uses monocular input with every 3 consecutive frames per camera. Waymo uses 3 frames (stride 4) from FRONT, SIDE\_LEFT, and SIDE\_RIGHT cameras, totaling 9 images per group. }
{\footnotesize\setlength{\tabcolsep}{.8mm}
{\begin{tabular}{lcccccccccccccccccc}
\toprule[1.2pt]                      
\multirow{3}{*}{\textbf{Methods}} & \multicolumn{6}{c}{\textbf{{KITTI (monocular)}}}   & \multicolumn{6}{c}{\textbf{{Waymo (3 cameras)}}}
\\ 
\cmidrule(lr){2-7} \cmidrule(lr){8-13} \cmidrule(lr){14-19}
& \multicolumn{3}{c}{\textbf{{Mean}}} & \multicolumn{3}{c}{\textbf{{Med.}}} &
\multicolumn{3}{c}{\textbf{{Mean}}} & \multicolumn{3}{c}{\textbf{{Med.}}} 
\\ 
\cmidrule(lr){2-4} \cmidrule(lr){5-7} \cmidrule(lr){8-10} \cmidrule(lr){11-13}
          & $ \text{Acc.} \downarrow $     & 
          $ \text{Comp.} \downarrow  $     &
          $ \text{NC} \uparrow  $     & 
          $ \text{Acc.} \downarrow   $     & 
          $ \text{Comp.} \downarrow    $     & 
          $ \text{NC} \uparrow   $     & 
          $ \text{Acc.} \downarrow   $     & 
          $ \text{Comp.} \downarrow    $     &
          $ \text{NC} \uparrow  $     & 
          $ \text{Acc.} \downarrow $     & 
          $ \text{Comp.} \downarrow  $     &
          $ \text{NC} \uparrow  $     &   
\\ 
\midrule

VGGT\cite{wang2025vggt}            & 1.489 & 0.690 & 0.918 & 1.329 & \textbf{0.535} & \textbf{0.971}    & 4.635 & 2.667 & 0.561 & 2.634 & 1.734 & 0.590  \\
StreamVGGT\cite{zhuo2025streaming}      & 1.078 & \textbf{0.495} & 0.899 & 0.867 & 0.390 & 0.949    & 4.598 & 2.626 & 0.564 & 2.567 & 1.789 & 0.592  \\
DynamicVGGT     &\textbf{0.901} &\underline{0.584} &\textbf{0.939} &\textbf{0.733} &\underline{0.464} &\underline{0.963}   &\textbf{4.021} &\textbf{2.390} &\textbf{0.562} &\textbf{1.971} &\underline{1.564} &\underline{0.603}                \\

\bottomrule[1.2pt]
\end{tabular}
}
}
\label{tab:point-estimation}
\end{table*}

Finally, we use scene-flow supervision for explicit Gaussian motion learning:
$
\mathcal{L}_{\mathrm{flow}} = \mathrm{MSE}(\mathbf{s}_{v,t}, \hat{\mathbf{s}}_{v,t}),
$
where $\mathbf{s}_{v,t}$ and $\hat{\mathbf{s}}_{v,t}$ denote the ground-truth and predicted scene flow, respectively. 
Compared with $\mathcal{L}_{\mathrm{temp}}$, which supervises coarse point displacement in the DPM space, $\mathcal{L}_{\mathrm{flow}}$ explicitly constrains motion in the Gaussian representation. 
The two losses therefore operate at different levels and are complementary rather than redundant.

